\bookchapter{Computational Interfaces Across the Three Volumes}\label{v3-connections-final}

The expanded computational record gives a concrete meaning to the handoff between the volumes. Volume I supplies substrate fields, propagation, clocks and physical records. Volume II supplies finite material sources, preparations, interactions and their receiving observables. Volume III supplies exact representations, executable transformations and retained results that can be passed to the next calculation. The interfaces are mathematical objects with declared maps, so a useful result can travel without changing what its symbols count or measure.\samref{V3C01,V3C02,V3C03}

\subsection{What crosses an interface}

A sufficient computational handoff can be written schematically as
\begin{equation}
 \mathcal B=(\mathcal S,\mu,\mathcal O,\mathcal C,\mathcal R,\mathcal D).
\end{equation}
Here $\mathcal S$ identifies the source and preparation, $\mu$ its state law or metric, $\mathcal O$ the operator, $\mathcal C$ its clock convention, $\mathcal R$ the requested receiver and $\mathcal D$ the retained dependencies. This tuple is an editorial organization of the interfaces developed in the preceding chapters. The particular calculation supplies its actual entries; it is not an additional physical law.

For a joint spin table, $\mathcal S$ is the signed integer Hamiltonian with ordered ports, $\mu$ is the labelled configuration count and $\mathcal R$ might be a rational-field DOS or a crossing query. There need be no dynamical clock in that static request. For a Gibbs response, $\mu$ is the original Gibbs inner product, $\mathcal O=N(I-P)$ and the retained history scale is $N/\xi$ uniform site attempts. For a signed-Write trajectory, the source path produces calibrated elapsed intervals received by the full carrier dynamics. Chapters~\ref{v3-exact-joint} and~\ref{v3-response-composition} give the corresponding implemented constructions.\samref{V3S01,V3S13,V3S14,V3S09}

This organization also explains how two calculations can share a native matrix operation while retaining different meanings. A Gram matrix may count states, measure a Gibbs residual, retain a completed arithmetic form or encode a receiver's observable directions. Its entries and metric establish the correspondence. Matching dimensions alone do not supply it.

\subsection{Volume I: source time, retained records and actual readouts}

The signed-Write integral supplies an explicit source-to-clock map. It receives the actual phase state and path, retains orientation and waiting, and sends a certified duration interval into the unchanged receiver. Different source orders with equal endpoints then produce different complete histories. Native dependency preservation carries that distinction through fresh sessions and later readouts.\samref{V3S08,V3S09}

The common-receiver and directional reconstruction work develops the complementary direction: determine what history can be recovered from a declared observation set. A receiver with signed phase and axis occupancy resolves a larger class than either partial observation. A source construction and its inverse therefore meet through the actual measured or computed record, with preparation, direction and sampling specified by the task. The application chapters retain those complete observation contracts.

The prime-event construction supplies another direct clock interface. Its microscopic Ising gate accumulates all-spin and output-spin activity along retained source paths. The exact output quotient has a simple two-state generator in its own activity coordinate, while actual sampling occurs at source-dependent stopping times. Chapter~\ref{v3-prime-event} develops how the hidden state and those stopping times remain observable. That source clock, the integer candidate index and a physical seconds assignment retain separate roles.\samref{V3C04}

The resulting computational contribution to Volume I is specific: an identified source can produce a duration, ordered record or response function that a field or clock construction receives through a declared map. This lets accumulation, propagation and clock questions use the completed computational operands directly while keeping their own physical definitions.\samref{V3C01}

\subsection{Volume II: exact phase sources and complete material receivers}

The phase-to-binary chart gives an immediate source correspondence. Its energy factor and aligned coordinate make the joint Ising table an executable phase-response object. Ordered boundaries retain which contact receives which sign. Conditional spectra and source constants remain available for thermal assembly; a ground-matched auxiliary interaction retains its own thermal factor.\samref{V3S02,V3S03,V3C02}

The later response work adds a second interface. Complete retained and hidden source functions generate Grams, Dirichlet forms and action residuals. A receiver carries the full signed mixed entries and the background terms that belong to that physical source. The latest edge calculation demonstrates a new source observable entering a sharper lower response through an explicit Schur extension. The source, clock and earlier maxima remain identifiable throughout.\samref{V3S14,V3S15}

The applied matter chapters show the same discipline in actual configurations: retain circulation ties, phase signs, incoming preparation and post-contact memory, then ask the declared readout question. Exact source selection, distinct-report selection and physical energy comparison are different operations. Each has its completed result and actual receiving interface. This gives Volume II reusable computational methods for source construction, history and response, with the material inventory preserved.\samref{V3C02}

\subsection{Volume III: several useful kinds of exact object}

The volume now develops several related return types. A full integer DOS retains the number of configurations at each energy. A joint table adds magnetization and ordered boundary information. A layered ground certificate retains the minimum, its exact degeneracy, a witness and exhaustive accounting. A certified matrix operation returns enclosures and residual bounds. A retained method trace supplies its operation/dependency graph. A prime or factor certificate answers its specified arithmetic question. Their shared computational infrastructure makes them composable when a receiving map is supplied.\samref{V3S01,V3S05,V3S06}

This distinction is productive when choosing a method. A thermal response needs the energy distribution or a justified source integration. A ground-energy question can use a minimum-preserving reduction. An inverse-history question needs the observations and ambiguity classes left by the actual receiver. An arithmetic transition estimate needs its signed source and completed complementary terms. A source-clock calculation needs the original path and clock map. The complete object requested by the next step determines which retained result can serve it.

\subsection{The record carried into the next three-volume pass}

The manuscript preserves the original twelve-chapter foundation and develops later completed work in separate chapters with an exact source register. Historical revisions remain attached to their results. Later installations, qualified methods and completed calculations have their own dated sources. Proposed successors and stopped or partial campaigns retain those states rather than being folded into completed totals.

That organization provides a usable starting point for the planned review of all three volumes. The next pass can follow a concrete chain: identify a receiving question, locate its existing source operands, determine the missing map or calculation, and select an applicable exact method. The present expansion supplies the completed work and its interfaces; the comparative selection of new research targets belongs to that subsequent review.


\subsection{Completed objects remain available to the next calculation}

The N5000 construction retains 1,530,259,992 primary joint records across six complete signed graphs, with exact integer multiplicities for every energy, magnetization and retained boundary in the specified sources. The N2000 physical analysis uses its own exact counts to derive residual entropy, excitation gaps and controlled thermal response. The prime-event program turns a source-level diagnosis into a prospective gate change and measures the improvement from one to forty-seven exact billion-index trajectories out of sixty-four per profile. Its later retained path reaches 1,039,515,000,000 decisions without an incorrect event before its intentional stop. Exact exchange reaches length thirteen, accounts for 67,108,864 states and composes with thirty-two local-matter families using the completed factors.

These results share a practical mathematical consequence. The completed calculation preserves the operands needed by a later operation: conditional integer counts for a field query, retained history for an inverse, source factors for a new matter spectrum, or the hidden-state law for a clock-sensitive response. Their constructions and validations establish the finite results presented here. Their explicit interfaces make those results reusable across the three volumes.
